\chapter{Компонентные observables рождения: selection rules без скорости}
\label{ch:v6-spectral-transition-componentwise-creation-observable-gate}

\section{Вопрос после отказа от единого пакета рождения}

Предыдущая архитектурная развилка сохранила класс пятнадцать как ledger
поколения и запретила читать его как multiplicity одного события. Теперь
нужно определить, какие компонентные утверждения остаются физически
проверяемыми без новой модели.

Сопоставляются четыре независимые структуры:
\begin{enumerate}
 \item тёплицевы ранги и веса компонент;
 \item физический сфалеронный спектральный поток;
 \item аномальные изменения зарядов;
 \item хиггс-разрешённые ранги локальных опор.
\end{enumerate}
Цель --- отделить точные selection rules от вероятностей, масс и
branching ratios, которых топология сама не задаёт.

\section{Классификационные компонентные числа}

Тёплицева граница раскладывается как
\begin{equation}
 \kappa_{15}=\kappa_{12}+\kappa_3,
\end{equation}
а нормированные следовые веса равны
\begin{equation}
 w_q=\frac{12}{105}=\frac4{35},
 \qquad
 w_\ell=\frac3{105}=\frac1{35}.
 \label{eq:v6-component-observable-toeplitz-weights}
\end{equation}
Их отношение равно
\begin{equation}
 w_q:w_\ell=4:1.
\end{equation}
Это точное утверждение о коэффициентном цикле и следовой нормировке. Но
в текущем родителе нет теоремы, превращающей $w_q,w_\ell$ в вероятности
создания кваркового и лептонного событий.

\section{Физический поток и зарядовые правила}

Для стандартного единичного сфалеронного перехода компонентные потоки
равны
\begin{equation}
 \operatorname{SF}_q:\operatorname{SF}_\ell=3:1.
 \label{eq:v6-component-observable-flow-ratio}
\end{equation}
При барионном весе $1/3$ каждой цветовой кварковой копии это даёт
\begin{equation}
 \Delta B=1,
 \qquad
 \Delta L=1,
 \qquad
 \Delta(B-L)=0.
 \label{eq:v6-component-observable-charge-rule}
\end{equation}
Связь топологического перехода, спектрального потока и аномальных
зарядов является стандартной
\cite{component-observable-klinkhamer-rupp2003,
component-observable-krs1985}.

Однако зарядовое отношение равно $1:1$, а не $3:1$ и не $4:1$.
Причина прозрачна: поток считает дублетные копии, заряд включает
физические веса, а тёплицев класс считает размер коэффициентного
носителя. Это разные observables.

\section{Хиггс-разрешённые опоры}

При $H\ne0$ конечный носитель дополнительно распадается как
\begin{equation}
 15=6_{\rm up}+6_{\rm down}+2_e+1_\nu.
 \label{eq:v6-component-observable-higgs-supports}
\end{equation}
Эти ранги определяют опоры блоков заряженного оператора и нейтринной
линии. Они также не являются числами частиц в одном переходе. В
частности, ранг нейтринного коварианта меняется $0\to1$ при прохождении
через $H=0$, но текущий функционал не задаёт частоту таких прохождений.

\section{Несовместимость единой вероятностной интерпретации}

После нормировки двухкомпонентные распределения, которые можно было бы
ошибочно принять за branching fractions, равны
\begin{align}
 p_{\rm rank}&=\left(\frac{12}{15},\frac3{15}\right)
 =\left(\frac45,\frac15\right),\\
 p_{\rm flow}&=\left(\frac34,\frac14\right),\\
 p_{B,L}&=\left(\frac12,\frac12\right).
\end{align}
Они попарно различны. В текущей архитектуре нет динамического принципа,
выбирающего одно из них или переводящего одно в другое.

\begin{theorem}[Selection rules не определяют branching ratios]
Текущий родитель точно фиксирует компонентные $K$-классы $(12,3)$,
сфалеронный поток $(3,1)$, зарядовые изменения
$(\Delta B,\Delta L)=(1,1)$ и хиггс-разрешённые ранги $(6,6,2,1)$.
Эти наборы являются различными спариваниями и не задают единой
вероятностной меры. Поэтому параметрически чистые branching ratio,
скорость рождения и абсолютная multiplicity из них не следуют.
\end{theorem}

\section{Почему скорость требует дополнительной динамики}

Даже в стандартной электрослабой теории сфалеронная скорость зависит от
температуры, барьерной энергии, калибровочных и хиггсовских параметров,
неравновесной эволюции и динамического prefactor. Она определяется
отдельным термическим или real-time расчётом, часто непертурбативным
\cite{component-observable-moore1998,
component-observable-barroso-moore2023}.

В проекте отсутствуют необходимые входы единой компонентной скорости:
\begin{enumerate}
 \item калиброванная энергия переходного барьера для проектного
 спектрального оператора;
 \item физическая температура или неравновесное распределение;
 \item real-time кинетический оператор и prefactor;
 \item мера флуктуаций, общая для компонент $12$ и $3$;
 \item нелинейное конечное состояние, интерпретируемое как наблюдаемая
 материя.
\end{enumerate}
Следовательно, формула вида $\Gamma\sim A e^{-E_*/T}$ не становится
предсказанием, пока $A,E_*,T$ не выведены одним родителем.

\section{Что остаётся строгим}

После настоящего гейта разрешены следующие утверждения:
\begin{itemize}
 \item классовый ledger: $(\kappa_{12},\kappa_3)$;
 \item физический поток стандартного сфалерона: $(3,1)$;
 \item аномальное правило: $\Delta B=\Delta L$, $\Delta(B-L)=0$;
 \item локальная возможность смены ранга $W_\nu:0\to1$ при $H=0$.
\end{itemize}
Не разрешены без новой динамики:
\begin{itemize}
 \item отношение числа рождаемых кварков и лептонов;
 \item абсолютная скорость события;
 \item масса или размер продукта;
 \item вероятность образования устойчивой материи;
 \item отождествление весов $4/35$ и $1/35$ с branching fractions.
\end{itemize}

\begin{nogocriterion}
Нельзя превращать ранги, индексы или нормированные следовые веса в
вероятности рождения без явно выведенной меры на траекториях и
real-time динамики. Совпадение суммы или красивое рациональное отношение
не заменяет расчёт скорости.
\end{nogocriterion}

\section{Следующий статусный гейт}

Следующий гейт
\path{version6_spectral_transition_dynamic_closure_status_gate}
должен свести полный ledger ветви спектрального рождения и ответить,
закрыта ли версия VI на уровне:
\begin{enumerate}
 \item кинематики и классификации;
 \item существования переходного седла;
 \item динамической вероятности;
 \item устойчивого конечного состояния;
 \item наблюдаемого физического предсказания.
\end{enumerate}

Нулевая гипотеза: классификационный и selection-rule уровни закрыты,
но динамическое рождение устойчивой наблюдаемой материи остаётся
незамкнутым.

Стоп-критерий: если отсутствуют и скорость, и устойчивый конечный
фермионный объект, дальнейшее продвижение внутри неизменённой
архитектуры фиксируется как круговое и требует новой модели.

\begin{protocol}
\begin{enumerate}
 \item компонентные классы: \textbf{$(12,3)$};
 \item нормированные тёплицевы веса: \textbf{$(4/35,1/35)$};
 \item физические потоки: \textbf{$(3,1)$};
 \item зарядовое правило: \textbf{$\Delta B=\Delta L=1$};
 \item сохранение $B-L$: \textbf{да};
 \item хиггс-разрешённые ранги: \textbf{$(6,6,2,1)$};
 \item единая компонентная вероятность: \textbf{не выведена};
 \item branching ratio: \textbf{не выведен};
 \item абсолютная скорость: \textbf{не выведена};
 \item масса и размер продукта: \textbf{не выведены};
 \item точные selection rules: \textbf{сохранены};
 \item динамическое замыкание: \textbf{не достигнуто}.
\end{enumerate}
\end{protocol}

Машинный сертификат сохранён в
\path{s2t_v6_spectral_transition_componentwise_creation_observable_gate_results.json}.

\begin{thebibliography}{9}
\bibitem{component-observable-klinkhamer-rupp2003}
F.~R.~Klinkhamer, C.~Rupp,
\emph{Sphalerons, Spectral Flow, and Anomalies},
Journal of Mathematical Physics 44 (2003), 3619--3639;
arXiv:hep-th/0304167.

\bibitem{component-observable-krs1985}
V.~A.~Kuzmin, V.~A.~Rubakov, M.~E.~Shaposhnikov,
\emph{On the Anomalous Electroweak Baryon Number Nonconservation in the
Early Universe},
Physics Letters B 155 (1985), 36--42.

\bibitem{component-observable-moore1998}
G.~D.~Moore,
\emph{Measuring the Broken Phase Sphaleron Rate Nonperturbatively},
Physical Review D 59 (1999), 014503; arXiv:hep-ph/9805264.

\bibitem{component-observable-barroso-moore2023}
M.~Barroso Mancha, G.~D.~Moore,
\emph{The Sphaleron Rate from 4D Euclidean Lattices},
Journal of High Energy Physics 01 (2023), 155;
arXiv:2210.05507.
\end{thebibliography}